\subsection{Quotient Rings and Maximal Ideals}

\begin{definition} \label{def:quotient-ring} \hlt{Quotient Rings}\\
Let $R$ be a ring and let $\mathcal{I}\trianglelefteq R$ be a two-sided ideal.
\begin{enumerate}[label=(\roman*)]
\setlength{\itemsep}{0pt}
\item The underlying set is the additive quotient group $R/\mathcal{I}=\{a+\mathcal{I}:a\in R\}$.
\item For $a,b\in R$, addition and multiplication are $(a+\mathcal{I})+(b+\mathcal{I})=(a+b)+\mathcal{I}$ and $(a+\mathcal{I})(b+\mathcal{I})=ab+\mathcal{I}$.
\item With these operations, $R/\mathcal{I}$ is the \hlt{quotient ring} of $R$ by $\mathcal{I}$.
\end{enumerate}
\end{definition}

\begin{proposition} \label{prop:quotient-ring-well-defined-by-ideal} \hlt{Ideals Make Quotient Multiplication Well Defined}\\
Let $R$ be a ring and let $\mathcal{I}\leq (R,+)$. The rule $(a+\mathcal{I})(b+\mathcal{I})\defeq ab+\mathcal{I}$ for $a,b\in R$ is well defined on $R/\mathcal{I}$ iff $\mathcal{I}\trianglelefteq R$.
\begin{proof}
\begin{enumerate}[label=\roman*.]
\setlength{\itemsep}{0pt}
\item[$(\Rightarrow)$] If the rule is well defined and $i\in\mathcal{I}$, then $i+\mathcal{I}=0_R+\mathcal{I}$.\\
For $r\in R$, $ri+\mathcal{I}=r0_R+\mathcal{I}$, and $ir+\mathcal{I}=0_Rr+\mathcal{I}$, so $ri,ir\in\mathcal{I}$. Thus $\mathcal{I}\trianglelefteq R$.
\item[$(\Leftarrow)$] If $\mathcal{I}\trianglelefteq R$ and $a',a,b',b\in R$ satisfy $a'=a+i$, $b'=b+j$ with $i,j\in\mathcal{I}$, then $a'b'-ab=aj+ib+ij\in\mathcal{I}$.\\
Hence $a'b'+\mathcal{I}=ab+\mathcal{I}$, so the product is independent of representatives.
\end{enumerate}
\hlt{Structural Commentary}: Ideals are exactly the additive subgroups compatible with multiplication of cosets.\\
\hlt{Key Insight}: Coset multiplication is well defined exactly for ideals.
\end{proof}
\end{proposition}

\begin{definition} \label{def:maximal-ideal} \hlt{Maximal Ideals}\\
Let $R$ be a ring and let $\mathcal{M}\trianglelefteq R$.
\begin{enumerate}[label=(\roman*)]
\setlength{\itemsep}{0pt}
\item $\mathcal{M}$ is \hlt{proper} if $\mathcal{M}\neq R$.
\item $\mathcal{M}$ is \hlt{maximal} if $\mathcal{M}\neq R$ and $(\forall\,\mathcal{J}\trianglelefteq R)\,[\mathcal{M}\subseteq\mathcal{J}\subseteq R\Rightarrow(\mathcal{J}=\mathcal{M}\lor\mathcal{J}=R)]$.
\end{enumerate}
\end{definition}

\begin{theorem} \label{thm:quotient-by-maximal-ideal-is-field} \hlt{Maximal Ideals and Fields}\\
Let $R$ be a commutative ring with identity and let $\mathcal{I}\trianglelefteq R$. Then $R/\mathcal{I}$ is a field $\iff$ $\mathcal{I}$ is maximal.
\begin{proof}
\begin{enumerate}[label=\roman*.]
\setlength{\itemsep}{0pt}
\item[$(\Rightarrow)$] Suppose $R/\mathcal{I}$ is a field. Then $1_R+\mathcal{I}\neq0_R+\mathcal{I}$, so $\mathcal{I}\neq R$. Let $\mathcal{I}\subsetneq\mathcal{J}\trianglelefteq R$ and choose $a\in\mathcal{J}\setminus\mathcal{I}$.\\
The nonzero coset $a+\mathcal{I}$ has an inverse $r+\mathcal{I}$, so $ra+\mathcal{I}=1_R+\mathcal{I}$. Thus $1_R-ra\in\mathcal{I}\subseteq\mathcal{J}$ and $ra\in\mathcal{J}$, whence $1_R\in\mathcal{J}$.\\
For every $x\in R$, $x=x1_R\in\mathcal{J}$, so $\mathcal{J}=R$ and $\mathcal{I}$ is maximal.
\item[$(\Leftarrow)$] Suppose $\mathcal{I}$ is maximal and let $a+\mathcal{I}\in R/\mathcal{I}$ satisfy $a+\mathcal{I}\neq0_R+\mathcal{I}$. Then $a\notin\mathcal{I}$.\\
Set $\mathcal{K}\defeq\langle\mathcal{I}\cup\{a\}\rangle$. Then $\mathcal{K}\trianglelefteq R$ and $\mathcal{I}\cup\{a\}\subseteq\mathcal{K}$ by Proposition \ref{prop:generated-ideal-smallest-containing-set}, so $\mathcal{I}\subsetneq\mathcal{K}$. Maximality forces $\mathcal{K}=R$.\\
By Definition \ref{def:ideal-and-generated-ideal}, an expression for $1_R\in\mathcal{K}$ can be regrouped as $1_R=i+ra$ for some $i\in\mathcal{I}$ and $r\in R$.\\
Hence $(r+\mathcal{I})(a+\mathcal{I})=1_R+\mathcal{I}$. The quotient ring is commutative with identity, and $\mathcal{I}\neq R$ because it is maximal.\\
Thus its zero and identity cosets are distinct, while every nonzero coset is invertible; hence $R/\mathcal{I}$ is a field.
\end{enumerate}
\hlt{Structural Commentary}: Maximal ideals are exactly the proper ideals whose quotient rings are fields.\\
\hlt{Key Insight}: Maximality makes every nonzero coset invertible.
\end{proof}
\end{theorem}

\begin{theorem} \label{thm:maximal-ideals-exist} \hlt{Existence of Maximal Ideals}\\
If $R$ is a commutative ring with identity and $1_R\neq0_R$, then $R$ has a maximal ideal.
\begin{proof}
Let $\mathcal{P}\defeq\{\mathcal{I}\trianglelefteq R:\mathcal{I}\neq R\}$, ordered by inclusion. Then $\mathcal{P}$ is a poset by Proposition \ref{prop:inclusion-and-reverse-inclusion-partial-orders}.\\
Also, $\{0_R\}=\langle\emptyset\rangle\trianglelefteq R$ by Proposition \ref{prop:generated-ideal-smallest-containing-set}; as $1_R\neq0_R$, $\{0_R\}\in\mathcal{P}$. Let $\mathcal{C}\subseteq\mathcal{P}$ be a chain. If $\mathcal{C}=\emptyset$, then $\{0_R\}$ upper bound for $\mathcal{C}$ in $\mathcal{P}$.\\
Assume $\mathcal{C}\neq\emptyset$ and set $\mathcal{U}\defeq\bigcup_{\mathcal{I}\in\mathcal{C}}\mathcal{I}$. Then $\mathcal{U}\trianglelefteq R$ by Proposition \ref{prop:intersections-and-chain-unions-of-ideals}.\\
Moreover, $\mathcal{U}\neq R$: if $1_R\in\mathcal{U}$, then $1_R\in\mathcal{I}$ for some $\mathcal{I}\in\mathcal{C}$, and $x=x1_R\in\mathcal{I}$ for every $x\in R$, contradicting $\mathcal{I}\in\mathcal{P}$.\\
Thus $\mathcal{U}\in\mathcal{P}$ is upper bound for $\mathcal{C}$. Every chain in $\mathcal{P}$ hence has an upper bound in $\mathcal{P}$, so $\mathcal{P}$ has a maximal element by Principle \ref{princ:zorns-lemma}.\\
By Definition \ref{def:maximal-ideal}, this element is a maximal ideal of $R$.\\
\hlt{Structural Commentary}: Maximal ideals are produced by maximality, not by an explicit formula in general.\\
\hlt{Key Insight}: Chain unions keep properness because containing $1$ is a finite obstruction.
\end{proof}
\end{theorem}
